% ====================================================================
% Файл: tasks/01_mechanics/kim04.tex
% Тема: Задание №4 (Статика, гидростатика, механические колебания и волны)
% ====================================================================

% === ЗАДАЧА 1 ===
% Тег: #МЕХАНИКА #КИМ_04 #ВАРИАНТ_1
\begin{taskbasic}[code={\label{task:dem26_v1_n4}}]
\textbf{Задача 1.} Школьник выполнял лабораторную работу по исследованию условий равновесия лёгкого рычага под действием двух сил. Полученные результаты он занёс в таблицу:
\begin{center}
\begin{tabular}{|c|c|c|c|}
\hline
$F_{1}$, Н & $l_{1}$, м & $F_{2}$, Н & $l_{2}$, м \\ \hline
40 & 0,6 & 120 & \\ \hline
\end{tabular}
\end{center}
Каково плечо $l_{2}$, если рычаг находится в равновесии?

\kim{4}
\end{taskbasic}
\answer{0{,}2~м}

% === ЗАДАЧА 2 ===
% Тег: #МЕХАНИКА #КИМ_04 #ВАРИАНТ_2
\begin{taskbasic}[code={\label{task:dem26_v2_n4}}]
\textbf{Задача 2.} Школьник выполнял лабораторную работу по исследованию условий равновесия лёгкого рычага под действием двух сил. Полученные результаты он занёс в таблицу:
\begin{center}
\begin{tabular}{|c|c|c|c|}
\hline
$F_{1}$, Н & $l_{1}$, м & $F_{2}$, Н & $l_{2}$, м \\ \hline
80 & & 120 & 0,2 \\ \hline
\end{tabular}
\end{center}
Каково плечо $l_{1}$, если рычаг находится в равновесии?

\kim{4}
\end{taskbasic}
\answer{0{,}3~м}

% === ЗАДАЧА 3 ===
% Тег: #МЕХАНИКА #КИМ_04 #ВАРИАНТ_3
\begin{taskbasic}[code={\label{task:dem26_v3_n4}}]
\textbf{Задача 3.} Период свободных малых колебаний математического маятника равен $0{,}8$~с. Каким станет период колебаний, если массу груза уменьшить в $4$ раза?

\kim{4}
\end{taskbasic}
\answer{0{,}8~с}

% === ЗАДАЧА 4 ===
% Тег: #МЕХАНИКА #КИМ_04 #ВАРИАНТ_4
\begin{taskbasic}[code={\label{task:dem26_v4_n4}}]
\textbf{Задача 4.} Период свободных малых колебаний математического маятника равен $0{,}8$~с. Какой станет частота колебаний, если массу груза увеличить в $4$ раза?

\kim{4}
\end{taskbasic}
\answer{1{,}25~Гц}

% === ЗАДАЧА 5 ===
% Тег: #МЕХАНИКА #КИМ_04 #ВАРИАНТ_5
\begin{taskbasic}[code={\label{task:dem26_v5_n4}}]
\textbf{Задача 5.} На рисунке показан график зависимости координаты $x$ от времени $t$ для одной из точек колеблющейся струны. Определите частоту колебаний струны.
\begin{center}
\begin{tikzpicture}[x=1.2cm, y=10cm]
    \draw[->] (0,-0.25) -- (0,0.25) node[left] {$x$, см};
    \draw[->] (0,0) -- (5.5,0) node[below right] {$t$, $10^{-3}$~с};
    \draw[dashed, gray] (0,0.2) -- (5,0.2);
    \draw[dashed, gray] (0,-0.2) -- (5,-0.2);
    \foreach \x [evaluate=\x as \labelx using int(\x*2)] in {1,2,3,4,5} {
        \draw (\x,0.02) -- (\x,-0.02);
        \node[below left] at (\x,0) {\labelx};
    }
    \node[left] at (0,0.2) {0,2};
    \node[left] at (0,0.1) {0,1};
    \node[left] at (0,-0.1) {-0,1};
    \node[left] at (0,-0.2) {-0,2};
    \node[below left] at (0,0) {0};
    \draw[ultra thick, domain=0:5, samples=100] plot (\x, {0.2*sin(180*\x)});
\end{tikzpicture}
\end{center}

\kim{4}
\end{taskbasic}
\answer{250~Гц}

% === ЗАДАЧА 6 ===
% Тег: #МЕХАНИКА #КИМ_04 #ВАРИАНТ_6
\begin{taskbasic}[code={\label{task:dem26_v6_n4}}]
\textbf{Задача 6.} На рисунке показан график зависимости координаты $x$ от времени $t$ для одной из точек колеблющейся струны. Какой путь проходит точка струны за два периода колебаний?
\begin{center}
\begin{tikzpicture}[x=1.2cm, y=10cm]
    \draw[->] (0,-0.25) -- (0,0.25) node[left] {$x$, см};
    \draw[->] (0,0) -- (5.5,0) node[below right] {$t$, $10^{-3}$~с};
    \draw[dashed, gray] (0,0.2) -- (5,0.2);
    \draw[dashed, gray] (0,-0.2) -- (5,-0.2);
    \foreach \x [evaluate=\x as \labelx using int(\x*2)] in {1,2,3,4,5} {
        \draw (\x,0.02) -- (\x,-0.02);
        \node[below] at (\x,0) {\labelx};
    }
    \node[left] at (0,0.2) {0,2};
    \node[left] at (0,0.1) {0,1};
    \node[left] at (0,-0.1) {-0,1};
    \node[left] at (0,-0.2) {-0,2};
    \node[below left] at (0,0) {0};
    \draw[ultra thick, domain=0:5, samples=100] plot (\x, {0.2*cos(180*\x)});
\end{tikzpicture}
\end{center}

\kim{4}
\end{taskbasic}
\answer{1{,}6~см}

% === ЗАДАЧА 7 ===
% Тег: #МЕХАНИКА #КИМ_04 #ВАРИАНТ_7
\begin{taskbasic}[code={\label{task:dem26_v7_n4}}]
\textbf{Задача 7.} Массивный шарик, подвешенный на лёгкой пружине, совершает свободные гармонические колебания вдоль вертикальной прямой. Во сколько раз нужно уменьшить массу груза, чтобы период его колебаний уменьшился в $4$ раза?

\kim{4}
\end{taskbasic}
\answer{16}

% === ЗАДАЧА 8 ===
% Тег: #МЕХАНИКА #КИМ_04 #ВАРИАНТ_8
\begin{taskbasic}[code={\label{task:dem26_v8_n4}}]
\textbf{Задача 8.} Массивный шарик, подвешенный на лёгкой пружине, совершает свободные гармонические колебания вдоль вертикальной прямой. Во сколько раз нужно увеличить массу груза, чтобы частота его колебаний уменьшилась в $1{,}5$ раза?

\kim{4}
\end{taskbasic}
\answer{2{,}25}

% === ЗАДАЧА 9 ===
% Тег: #МЕХАНИКА #КИМ_04 #ВАРИАНТ_9
\begin{taskbasic}[code={\label{task:dem26_v9_n4}}]
\textbf{Задача 9.} Скорость звука в алюминиевой трубе составляет $5$~км/с. Длина звуковой волны в трубе равна $2{,}5$~м. Какова частота колебаний источника звука?

\kim{4}
\end{taskbasic}
\answer{2000~Гц}

% === ЗАДАЧА 10 ===
% Тег: #МЕХАНИКА #КИМ_04 #ВАРИАНТ_10
\begin{taskbasic}[code={\label{task:dem26_v10_n4}}]
\textbf{Задача 10.} Скорость звука в мраморной колонне составляет $3{,}8$~км/с. Длина звуковой волны в колонне равна $1{,}9$~м. Каков период колебаний источника звука?

\kim{4}
\end{taskbasic}
\answer{0{,}0005~с}

% === ЗАДАЧА 11 ===
% Тег: #МЕХАНИКА #КИМ_04 #ВАРИАНТ_11
\begin{taskbasic}[code={\label{task:dem26_v11_n4}}]
\textbf{Задача 11.} Железный кубик плавает в ртути. Ребро кубика равно $10$~см. Определите силу Архимеда, действующую на кубик.

\kim{4}
\end{taskbasic}
\answer{78~Н}

% === ЗАДАЧА 12 ===
% Тег: #МЕХАНИКА #КИМ_04 #ВАРИАНТ_12
\begin{taskbasic}[code={\label{task:dem26_v12_n4}}]
\textbf{Задача 12.} Медный кубик, подвешенный на нити, полностью погружён в воду и не касается дна сосуда. Ребро кубика равно $3$~см. Определите силу Архимеда, действующую на кубик.

\kim{4}
\end{taskbasic}
\answer{0{,}27~Н}

% === ЗАДАЧА 13 ===
% Тег: #МЕХАНИКА #КИМ_04 #ВАРИАНТ_13
\begin{taskbasic}[code={\label{task:dem26_v13_n4}}]
\textbf{Задача 13.} Во сколько раз увеличится период малых свободных колебаний математического маятника, если длину нити увеличить в $9$ раз, а массу груза уменьшить в $2{,}25$ раза?

\kim{4}
\end{taskbasic}
\answer{3}

% === ЗАДАЧА 14 ===
% Тег: #МЕХАНИКА #КИМ_04 #ВАРИАНТ_14
\begin{taskbasic}[code={\label{task:dem26_v14_n4}}]
\textbf{Задача 14.} Во сколько раз уменьшится частота малых свободных колебаний математического маятника, если длину нити увеличить в $2{,}25$ раза, а массу груза уменьшить в $4$ раза?

\kim{4}
\end{taskbasic}
\answer{1{,}5}

% === ЗАДАЧА 15 ===
% Тег: #МЕХАНИКА #КИМ_04 #ВАРИАНТ_15
\begin{taskbasic}[code={\label{task:dem26_v15_n4}}]
\textbf{Задача 15.} Груз, подвешенный на лёгкой пружине жёсткостью $32$~Н/м, совершает свободные вертикальные гармонические колебания. Пружину какой жёсткости надо взять вместо этой пружины, чтобы период свободных вертикальных колебаний этого груза стал в $2{,}5$ раза меньше?

\kim{4}
\end{taskbasic}
\answer{200~Н/м}

% === ЗАДАЧА 16 ===
% Тег: #МЕХАНИКА #КИМ_04 #ВАРИАНТ_16
\begin{taskbasic}[code={\label{task:dem26_v16_n4}}]
\textbf{Задача 16.} Груз массой $400$~г, подвешенный на лёгкой пружине, совершает свободные вертикальные гармонические колебания. На сколько надо увеличить массу груза, чтобы частота его свободных вертикальных колебаний на этой пружине стала в $2$ раза меньше?

\kim{4}
\end{taskbasic}
\answer{1{,}2~кг}

% === ЗАДАЧА 17 ===
% Тег: #МЕХАНИКА #КИМ_04 #ВАРИАНТ_17
\begin{taskbasic}[code={\label{task:dem26_v17_n4}}]
\textbf{Задача 17.} Невесомый рычаг находится в равновесии. Модуль силы $F_{1} = 12$~Н, её плечо равно $80$~см, а модуль силы $F_{2} = 48$~Н. Определите длину рычага.

\kim{4}
\end{taskbasic}
\answer{100~см}

% === ЗАДАЧА 18 ===
% Тег: #МЕХАНИКА #КИМ_04 #ВАРИАНТ_18
\begin{taskbasic}[code={\label{task:dem26_v18_n4}}]
\textbf{Задача 18.} Невесомый рычаг находится в равновесии. Модуль силы $F_{2} = 12$~Н, её плечо равно $10$~\text{см}. Каков модуль силы $F_{1}$, если её плечо равно $15$~см?

\kim{4}
\end{taskbasic}
\answer{8~Н}

% === ЗАДАЧА 19 ===
% Тег: #МЕХАНИКА #КИМ_04 #ВАРИАНТ_19
\begin{taskbasic}[code={\label{task:dem26_v19_n4}}]
\textbf{Задача 19.} Ученик выполнял лабораторную работу по исследованию условий равновесия рычага под действием двух сил: $\vec{F}_{1}$ и $\vec{F}_{2}$, плечи которых равны соответственно $l_{1}$ и $l_{2}$. Он внёс результаты измерений в таблицу:
\begin{center}
\begin{tabular}{|c|c|c|c|}
\hline
$F_{1}$, Н & $l_{1}$, м & $F_{2}$, Н & $l_{2}$, м \\ \hline
& 0,4 & 20 & 1,2 \\ \hline
\end{tabular}
\end{center}
Определите модуль силы $\vec{F}_{1}$, если рычаг находится в равновесии. Массой рычага пренебречь.

\kim{4}
\end{taskbasic}
\answer{60~Н}

% === ЗАДАЧА 20 ===
% Тег: #МЕХАНИКА #КИМ_04 #ВАРИАНТ_20
\begin{taskbasic}[code={\label{task:dem26_v20_n4}}]
\textbf{Задача 20.} Ученик выполнял лабораторную работу по исследованию условий равновесия рычага под действием двух сил: $\vec{F}_{1}$ и $\vec{F}_{2}$, плечи которых равны соответственно $l_{1}$ и $l_{2}$. Он внёс результаты измерений в таблицу:
\begin{center}
\begin{tabular}{|c|c|c|c|}
\hline
$F_{1}$, Н & $l_{1}$, м & $F_{2}$, Н & $l_{2}$, м \\ \hline
80 & 0,25 & 20 & \\ \hline
\end{tabular}
\end{center}
Определите длину плеча $l_{2}$, если рычаг находится в равновесии. Массой рычага пренебречь.

\kim{4}
\end{taskbasic}
\answer{1~м}